% Common preamble for all lab reports
\documentclass[a4paper, 12pt]{article}

% Language and encoding
\usepackage[T2A]{fontenc}
\usepackage[utf8]{inputenc}
\usepackage[ukrainian]{babel}

% Typography
\usepackage{microtype}
\usepackage{lmodern}

% Page geometry
\usepackage[margin=25mm, headheight=14pt]{geometry}

% Math
\usepackage{amsmath, amssymb}

% Graphics and figures
\usepackage{graphicx}
\usepackage{float}
\usepackage{subcaption}
\usepackage[justification=centering]{caption}

% Tables
\usepackage{tabularx}
\usepackage{booktabs}
\usepackage{array}

% Colors and hyperlinks
\usepackage{xcolor}
\usepackage[colorlinks=true, linkcolor=blue!60!black, urlcolor=blue!60!black, citecolor=blue!60!black]{hyperref}

% Lists
\usepackage{enumitem}
\setlist{nosep, leftmargin=*}

% Section formatting
\usepackage{titlesec}
\titleformat{\section}{\large\bfseries}{\thesection.}{0.5em}{}
\titleformat{\subsection}{\normalsize\bfseries}{\thesubsection.}{0.5em}{}

% Header/footer
\usepackage{fancyhdr}
\pagestyle{fancy}
\fancyhf{}
\fancyhead[L]{\small\leftmark}
\fancyhead[R]{\small\thepage}
\renewcommand{\headrulewidth}{0.4pt}
\fancypagestyle{plain}{\fancyhf{}\fancyfoot[C]{\thepage}\renewcommand{\headrulewidth}{0pt}}

% Placeholder image command
\newcommand{\placeholder}[2][0.8\textwidth]{%
  \begin{figure}[H]
    \centering
    \fbox{\parbox{#1}{\vspace{4cm}\centering\textit{#2}\vspace{4cm}}}
    \caption{#2}
  \end{figure}
}

% Placeholder for inline result values
\newcommand{\result}[1]{\textbf{\textcolor{red!70!black}{[#1]}}}

\usepackage{../common/labstyle}

\labnum{2}
\labtitle{Гідродинаміка потоку}

\begin{document}
\maketitle
\tableofcontents
\newpage

% ===================================================================
\section{Вихідні дані}
% ===================================================================

\begin{table}[H]
  \centering
  \caption{Вихідні дані (варіант~13)}
  \begin{tabularx}{0.75\textwidth}{Xr}
    \toprule
    \textbf{Параметр} & \textbf{Значення} \\
    \midrule
    Об'єкт дослідження         & Конус (воронка) \\
    Положення відносно потоку  & Широким кінцем проти потоку \\
    Швидкість на вході (Inlet) & 1{,}3~м/с \\
    Діаметр труби              & 100~мм \\
    Довжина труби              & 300~мм \\
    Матеріал потоку            & Water \\
    Макс.\ габарити об'єкта    & $50 \times 50 \times 50$~мм \\
    \bottomrule
  \end{tabularx}
\end{table}

\subsection{Ключові формули}

Число Рейнольдса: $\mathrm{Re} = \rho U D / \mu$.
Рівняння Бернуллі: $p_1 + \tfrac{1}{2}\rho U_1^2 = p_2 + \tfrac{1}{2}\rho U_2^2$.
Модель турбулентності $k$-$\varepsilon$: турбулентна в'язкість $\mu_t = \rho C_\mu k^2/\varepsilon$.

% ===================================================================
\section{Послідовність виконання}
% ===================================================================

Робота виконана через ANSYS Workbench Journaling (\texttt{lab2.wbjn}) з солвером CFX.

\begin{lstlisting}[style=scheme]
1. CreateSystem("Fluid Flow (CFX)")
2. Import lab2.step (1 body)
3. Enter Meshing
   |-- Print face centroids (diagnostic)
   |-- NS "Inlet"  -> faces at X < 0.5mm
   |-- NS "Outlet" -> faces at X > 299.5mm
   |-- NS "Wall"   -> rest (pipe cylinder + cone void)
   |-- Inflation -> scoped to wall_faces, 5 layers, h1=0.3mm
   |-- Sizing 1.5mm -> scoped to cone_faces (100 < X < 200mm)
   |-- Mesh 3mm -> mesh.png
4. Enter CFX-Pre
   |-- Import CCL: FluidDomain=B1, Water, k-eps
   |   Inlet: 1.3 m/s, Medium Turbulence
   |   Outlet: 0 Pa avg static
   |   Wall: No Slip, Smooth
   |   Solver: High Res, 300 iter, RMS 1e-4
5. Solve (system.Update)
6. Enter CFX-Post
   |-- Import post CCL (all 7 objects at once)
   |-- Plane: XZ at Y=0
   |-- Export: velocity, pressure, ssr, tke, tef,
   |           streamline, vector
   |-- Function Calculator -> console log (9 calcs)
   '-- geometry.png
   NOTE: parse numeric values from solver .out log
\end{lstlisting}

\subsection{Геометрія (CadQuery $\to$ STEP)}

Скрипт \texttt{geometry/lab2.py}: циліндрична труба ($D=100$~мм, $L=300$~мм вздовж осі~$X$), конус ($50\times50\times50$~мм) по центру, Boolean Subtract $\to$ одне рідинне тіло. Експорт: \texttt{geometry/step/lab2.step}.

\placeholder{Геометрія труби з вирізаним конусом (воронкою)}

\placeholder{Розрахункова сітка (Mesh) з ущільненням навколо об'єкта}

% ===================================================================
\section{Результати розрахунку}
% ===================================================================

\subsection{Velocity}
\placeholder{Contour: розподіл швидкості (Velocity)}

\subsection{Pressure}
\placeholder{Contour: розподіл тиску (Pressure)}

\subsection{Shear Strain Rate}
\placeholder{Contour: Shear Strain Rate}

\subsection{Turbulence Kinetic Energy}
\placeholder{Contour: Turbulence Kinetic Energy}

\subsection{Turbulence Eddy Frequency}
\placeholder{Contour: Turbulence Eddy Frequency}

\subsection{Streamline}
\placeholder{Лінії току (Streamline)}

\subsection{Vector}
\placeholder{Вектори швидкості (Vector)}

\subsection{Function Calculator}

\begin{table}[H]
  \centering
  \caption{Результати Function Calculator}
  \begin{tabularx}{0.85\textwidth}{Xlrr}
    \toprule
    \textbf{Параметр} & \textbf{Функція} & \textbf{Зона} & \textbf{Значення} \\
    \midrule
    Velocity     & maxVal  & Domain & \result{---} \\
    Velocity     & areaAve & Inlet  & \result{---} \\
    Velocity     & areaAve & Outlet & \result{---} \\
    Pressure     & maxVal  & Domain & \result{---} \\
    Pressure     & minVal  & Domain & \result{---} \\
    Pressure     & areaAve & Inlet  & \result{---} \\
    Shear Strain Rate & maxVal & Domain & \result{---} \\
    Turb.\ Kinetic Energy & maxVal & Domain & \result{---} \\
    Turb.\ Eddy Frequency & maxVal & Domain & \result{---} \\
    \bottomrule
  \end{tabularx}
\end{table}

% ===================================================================
\section{Аналіз та висновок}
% ===================================================================

\result{Детальний опис процесів на кожній епюрі:
\begin{itemize}
  \item Зони застою (stagnation) --- перед широким кінцем конуса;
  \item Прискорення потоку --- у звуженні між конусом та стінками труби;
  \item Завихрення (recirculation) --- за конусом у кільвій зоні (wake);
  \item Розподіл TKE та Eddy Frequency;
  \item Загальний висновок про вплив конуса на гідродинамічний опір.
\end{itemize}}

\end{document}
